% Generalized way to enter data into a table and format the cell.
% Formats the cell height in a nice way.
% Usage: \wksentry{<cell width>}{<Text in cell>}
\newcommand{\wksentry}[2]{
        \begin{minipage}[c][\totalheight+.5cm][c]{#1}
        \raggedright
        #2
        \end{minipage}
}%end wksentry


% cwksentry() 
% 
% Another nice way to put data in a cell.  Centers the text
% in a minipage.
%  
% Usage: \cwksentry{<cell width>}{<Text in cell>}
\newcommand{\cwksentry}[2]{
  \begin{minipage}[c][\totalheight+.5cm][c]{#1}
    \begin{center}
      #2
    \end{center}
  \end{minipage}
} % end cwksentry

% Shade a cell in a table
\newcommand{\graycell}[0]{\cellcolor[gray]{.8}}

% Highlight the requirements for a certain step
\newcommand{\reqtext}[1]{
    \begin{center} 
        \textbf{#1}
    \end{center}
}%end reqtext

% A terminal window showing a command issued by root and the response
% Usage: \rootcmd{<command>}{<output>}
\newcommand{\rootcmd}[2]{
\vspace{-.5\baselineskip}
\begin{center}
    \singlespacing
    \vspace{-\baselineskip}
    \begin{boxedminipage}[t]{\textwidth}
        \texttt{\# #1} \\
        #2
    \end{boxedminipage}
\end{center}
}% end rootcmd

% A terminal window showing a command issued by an ordinary user
% and the response
% Usage: \usercmd{<command>}{<output>}
\newcommand{\usercmd}[2]{
\vspace{-.5\baselineskip}
\begin{center}
    \singlespacing
    \vspace{-\baselineskip}
    \begin{boxedminipage}[t]{\textwidth}
        \textbf{\$} \texttt{#1} \\
        #2
    \end{boxedminipage}
\end{center}
}% end usercmd

% boxcmd() 
%
% Typewriter formatted text inside a box.  Use this to draw up any
% kind of terminal communication with any kind of prompt
%
% Usage: \boxcmd{<text>}
\newcommand{\boxcmd}[1]{
\vspace{-.5\baselineskip}
\begin{center}
    \singlespacing
    \vspace{-\baselineskip}
    \begin{boxedminipage}[t]{\textwidth}
        \noindent
        \texttt{#1}
    \end{boxedminipage}
\end{center}
}% end \boxcmd

% boxtext()
% 
% Simply text inside of a box of width \textwidth.
%
% Usage: \boxtext{<text>}
\newcommand{\boxtext}[1]{
  \vspace{-.5\baselineskip}
  \begin{center}
    \singlespacing
    \vspace{-\baselineskip}
    \begin{boxedminipage}[t]{\textwidth}
      \noindent
      {#1}
    \end{boxedminipage}
  \end{center}
}% end \boxtext

% onecodeline()
%
% Uses \path from the url package to verbatim-print input all on one
% line.  This is handy if you have a long line of input for verbatim
% that you'd like to put on two lines in your editor, but not in the output.
\newcommand{\onecodeline}[1]{
  \vspace{-.5\baselineskip}
  \begin{center}
    \colorbox{codegray}{
      \begin{minipage}[c]{\textwidth - 2\tabcolsep}
        \path{#1}
      \end{minipage}
    }
  \end{center}
}


% filesnip()
%
% Show a snippet of a file's contents in a box with the filename as
% the title.  The filename is shaded.
%
% Requires: colortbl
%
% Usage: \filesnip{<filename>}{<snippet>}
\newcommand{\filesnip}[2]{
\begin{center}
    \vspace{-\baselineskip}
    \begin{tabular}{|l|} \hline
        \rowcolor[gray]{0.8}
        \begin{minipage}[c]{\textwidth - 2\tabcolsep}
            \textbf{File:}
            #1
        \end{minipage}\\
        \begin{minipage}[c]{\textwidth - 2\tabcolsep}
            \vspace{0.5\baselineskip}
            $\cdots$ \\
            \begin{minipage}[c]{\textwidth - 2\tabcolsep}
                #2
            \end{minipage}
            $\cdots$\\
            \vspace{-0.5\baselineskip}
        \end{minipage}\\
    \hline
    \end{tabular}
\end{center}
}

% Show the entire contents of a file (not a snippet) in a box with the
% filename as the title.  The filename is shaded.
% Usage: \filecont{<filename>}{<text>}
\newcommand{\filecont}[2]{
\begin{center}
    \vspace{-\baselineskip}
    \begin{tabular}{|l|} \hline
        \rowcolor[gray]{0.8}
        \begin{minipage}[c]{\textwidth - 2\tabcolsep}
            \textbf{File:}
            #1 (Entire file)
        \end{minipage}\\
        \begin{minipage}[c]{\textwidth - 2\tabcolsep}
        \singlespacing
        #2
        \vspace{1mm}
        \end{minipage}\\
    \hline
    \end{tabular}
\end{center}
}

% Make a note next to an icon.  Like, here's a stop sign and some
% notes about why what you're doing is bad
% Usage: \picnote{<image>}{<scale factor>}{<Text in note>}
\newcommand{\picnote}[3]{
    \noindent
    \begin{tabular}{lc@{}}
    \begin{minipage}[c][\height+1cm][c]{2cm}
        \includegraphics[clip,scale=#2]{#1}
    \end{minipage}
    &
    \begin{minipage}[c][\height+1cm][c]{\textwidth-2cm-3\tabcolsep}
    #3
    \end{minipage}
    \end{tabular}
}% end picnote{}

% Make a warning note. This means place some text next to a warning
% image like a stop sign
% Usage: \rednote{<Text in note>}
\newcommand{\rednote}[1]{
    \noindent \\
    \begin{tabular}{lc@{}}
    \begin{minipage}[c][\height+1cm][c]{2cm}
        \includegraphics[clip,scale=1.5]{doctools/latex/icons/error.eps}
    \end{minipage}
    &
    \begin{minipage}[c][\height+1cm][c]{\textwidth-2cm-3\tabcolsep}
    #1
    \end{minipage}
    \end{tabular}
}

% Make a general note.  This means place some text next to a little
% sticky note icon.
% Usage: \nailnote{<Text in note>}
\newcommand{\nailnote}[1]{
    \noindent \\
    \begin{tabular}{lc@{}}
    \begin{minipage}[c][\height+1cm][c]{2cm}
        \includegraphics[clip,scale=0.6]{doctools/latex/icons/notes.eps}
    \end{minipage}
    &
    \begin{minipage}[c][\height+1cm][c]{\textwidth-2cm-3\tabcolsep}
    #1
    \end{minipage}
    \end{tabular}
}

% Make a note with a question mark next to it.
\newcommand{\faqnote}[1]{
    \noindent \\
    \begin{tabular}{lc@{}}
    \begin{minipage}[c][\height+1cm][c]{2cm}
        \includegraphics[clip,scale=1.5]{doctools/latex/icons/question.eps}
    \end{minipage}
    &
    \begin{minipage}[c][\height+1cm][c]{\textwidth-2cm-3\tabcolsep}
    #1
    \end{minipage}
    \end{tabular}
}

% Make a note indicating that this function can be enabled with a PC
% Puts a picture reminding somebody of a pc next to some text.
% Usage: \pcnote{<Text in note>}
\newcommand{\pcnote}[1]{
    \noindent \\
    \begin{tabular}{lc}
    \begin{minipage}[c][\height+1cm][c]{2cm}
        \includegraphics[clip,scale=0.8]{doctools/latex/icons/terminal.eps}
    \end{minipage}
    &
    \begin{minipage}[c][\height+1cm][c]{\textwidth-2.5cm}
    #1
    \end{minipage}
    \end{tabular}
}

% Add a filename to a caption if the document is marked as a draft
% Need to add the line: \newcommand{\isdraft}{1} in the main document
% Also, if using this in a caption, use the full caption form:
% \caption[short caption]{long caption}
% Usage: \draftnote{<text>}
\newcommand{\draftnote}[1]{
    \ifthenelse{\equal{\isdraft}{1}}{\textcolor{green}{\\ \textsl{#1}}}{}
}

% An ascii command definition.  Just places the command in typewriter
% font on one side of the page, and its description on the other.
% Usage: \asciicmd{<command>}{<description}
%
% Very important: if you want to add a label or index entry for the
% command, put the \label or \index command inside the {<description>}
% argument.  Otherwise the \label or \index commands may end up on the
% page preceeding the minipage -- the index refers to the ascii
% command on the page before it actually appears.
\newcommand{\asciicmd}[2]{
    \noindent
    \begin{tabular}{@{}l@{}c@{}}
    \begin{minipage}[c][\totalheight+1cm][c]{\cmdwidth}
        \texttt{#1}
    \end{minipage}
    &
    \begin{minipage}[c][\totalheight+1cm][c]{\textwidth-\cmdwidth}
    #2
    \end{minipage}
    \end{tabular}
}

% Make a figure sitting against the left margin that text wraps around.
%
% The number of shortened text lines produced must be specified to avoid
% conflicts with floats below the figure.  Change this number if your
% floats don't move over to the left margin.
%
% Usage: \sidewrap{<image>}{<scale factor>}{<number of short lines>}
% Must define \sidefigwidth somewhere.  A nice value is 6cm
\newcommand{\sidewrap}[3]{
    \begin{wrapfigure}[#3]{L}[0cm]{\sidefigwidth}
        \begin{minipage}[c][]{\sidefigwidth}
            \begin{center}
                \includegraphics[clip,scale=#2]{#1}
            \end{center}
        \end{minipage}
    \end{wrapfigure}
}

% Make two pieces of text -- one up against the left margin, and the
% other against the right.  Like in a resume.
\newcommand{\leftright}[2]{
    \begin{tabular}{@{}l@{}r@{}}
        \begin{minipage}[t][\totalheight]{0.75\textwidth}
            {\noindent #1}
        \end{minipage}
        &\begin{minipage}[t][\totalheight]{0.25\textwidth}
            \begin{flushright}
            {\noindent #2}
            \end{flushright}
        \end{minipage}\\
    \end{tabular}
}

% Make a centered piece of text between two horizontal rules
\newcommand{\sandwich}[1]{
    \rule{\textwidth}{0.4pt}
    \begin{minipage}[c][\totalheight]{\textwidth}
        \begin{center}
            {\noindent #1}
        \end{center}
    \end{minipage}
    \vspace{-1.5\baselineskip}\\
    \rule{\textwidth}{0.4pt}
}

% Make a horizontal line dividing two sections of text
\newcommand{\hdivline}{
    \begin{center}
        \vspace{-0.5\baselineskip}
        \rule{\textwidth}{1pt}
    \end{center}
}

% codecomment
% Format text into a comment for a code snippet
\newcommand{\codecomment}[1]{
    \textsl{\textcolor{green}{{#1}}}
}

% Format the text describing a key on the keyboard
\newcommand{\kbkey}[1]{
    $<$\textit{#1}$>$
}

% Format the text describing a GUI's menu item
\newcommand{\menuitem}[1]{
	\textbf{#1}
}

% buttname(name of button)
% Format text into the name of a front panel button
\newcommand{\buttname}[1]{
    \texttt{[\MakeUppercase{#1}]}
}

% bncname(name of bnc)
% Format text into the name of a panel BNC connector
\newcommand{\bncname}[1]{
    \texttt{\MakeUppercase{#1}}
}

% progname(name of program)
% Format text into the name of an external program, like gnuplot
% or gnetlist.
\newcommand{\progname}[1]
{
    \textbf{{#1}}
}

% urlname(name of a url or file)
% Format text into the name of a url or file.
\newcommand{\urlname}[1]
{
    \textit{{#1}}
}

% mfgpart(name of a part)
% Format text into the name of a part, like the opa134 opamp or the
% atmega169 uC.
\newcommand{\mfgpart}[1]
{
    \MakeUppercase{#1}
}



